\documentclass{article}
\usepackage[a4paper,portrait,hmargin=10.5mm,top=8mm,bottom=8mm]{geometry}
\usepackage[T1]{fontenc}
\usepackage{lmodern}
\usepackage{amsmath,amssymb}
% Use scalable mathematical operators at the reference text sizes.
\input{omxlmex.fd}
\DeclareFontShape{OMX}{lmex}{m}{n}{<->lmex10}{}
\usepackage{xcolor,colortbl}
\usepackage{tabularx,array}
\pagestyle{empty}
\definecolor{navy}{HTML}{172C3C}
\definecolor{coral}{HTML}{BF604D}
\definecolor{rulecolor}{HTML}{A5AFB6}
\color{navy}
\setlength{\parindent}{0pt}
\setlength{\parskip}{0pt}
\setlength{\arrayrulewidth}{0.3pt}
% Central typography controls, explicitly applied in table cells.
\newcommand{\TitleFont}{\fontsize{25}{27}\selectfont\bfseries}
\newcommand{\KickerFont}{\fontsize{10}{12}\selectfont\bfseries}
\newcommand{\HeadingFont}{\fontsize{12.5}{15}\selectfont\bfseries}
\newcommand{\DefinitionFont}{\fontsize{9}{10.3}\selectfont}
\newcommand{\SymbolFont}{\fontsize{10}{11}\selectfont}
\newcommand{\ExampleFont}{\fontsize{8.8}{10.2}\selectfont}
\newcommand{\TableHeadFont}{\fontsize{8.5}{10}\selectfont\bfseries}
\newcommand{\FooterFont}{\fontsize{8.5}{10}\selectfont}
\DeclareMathSizes{10}{10}{7}{5}
\DeclareMathSizes{9}{9}{6.5}{5}
\DeclareMathSizes{8.8}{8.8}{6.3}{5}

\newlength{\SymbolWidth}
\newlength{\MeaningWidth}
\newlength{\RowPadding}
\setlength{\RowPadding}{0.7mm}
% Measure the complete cell (including wrapped text and tall mathematics),
% then reserve the same clearance above and below its actual bounds.
\newsavebox{\CellContent}
\newcommand{\paddedcell}[1]{%
  \sbox{\CellContent}{\begin{minipage}[t]{\linewidth}\raggedright #1\end{minipage}}%
  \raisebox{0pt}[\dimexpr\ht\CellContent+\RowPadding\relax]%
    [\dimexpr\dp\CellContent+\RowPadding\relax]{\usebox{\CellContent}}%
}
% Center only complete symbol rows, using the actual bounds of tall mathematics.
% Headers and the letter reference retain their original padded-cell layout.
\newcommand{\centeredcell}[1]{%
  \sbox{\CellContent}{\begin{minipage}[t]{\linewidth}\raggedright #1\end{minipage}}%
  \raisebox{\dimexpr(\dp\CellContent-\ht\CellContent)/2\relax}%
    [\dimexpr(\ht\CellContent+\dp\CellContent)/2+\RowPadding\relax]%
    [\dimexpr(\ht\CellContent+\dp\CellContent)/2+\RowPadding\relax]%
    {\usebox{\CellContent}}%
}
% Text-style arrays align relations without enlarging fractions or operators.
% Each r/E pair holds a left-hand side and its relation/right-hand side.
% E supplies an empty math atom to retain normal spacing before a relation.
\newcommand{\ExampleLineStretch}{1.12}
\newcommand{\exmath}[2][r@{}E]{{\renewcommand{\arraystretch}{\ExampleLineStretch}%
  \ifcsname NC@find@E\endcsname\else\newcolumntype{E}{>{{}}l}\fi%
  \ensuremath{\begin{array}[t]{@{}#1@{}}#2\end{array}}}}
\newcommand{\symrow}[3]{%
  \centeredcell{\centering\SymbolFont\ensuremath{\displaystyle #1}}&%
  \centeredcell{\DefinitionFont #2}&\centeredcell{\ExampleFont #3}\\\hline
}
% Measure each example line so paired symbols and meanings share its baselines,
% including fractions and wrapped meanings. Spacing settings stay local.
\newsavebox{\PairFirst}
\newsavebox{\PairSecond}
\newcommand{\pairstrut}[1]{\rule[-\dp#1]{0pt}{\dimexpr\ht#1+\dp#1\relax}}
\newcommand{\pairlines}[3]{%
  \begin{tabular}[t]{@{}#1@{}}%
  \pairstrut{\PairFirst}#2\\%
  \pairstrut{\PairSecond}#3%
  \end{tabular}%
}
% Paired symbols, meanings, and examples; wide blocks can supply their columns.
\newcommand{\pairedrow}[6]{%
  \pairedrowwithcols{r@{}E@{\quad}r@{}E}{#1}{#2}{#3}{#4}{#5}{#6}}
\newcommand{\pairedrowwithcols}[7]{%
  \noalign{\begingroup
  \renewcommand{\ExampleLineStretch}{1}%
  \global\setbox\PairFirst=\hbox{\ExampleFont\exmath[#1]{#6}%
    \DefinitionFont\parbox[t]{\MeaningWidth}{\raggedright #4}}%
  \global\setbox\PairSecond=\hbox{\ExampleFont\exmath[#1]{#7}%
    \DefinitionFont\parbox[t]{\MeaningWidth}{\raggedright #5}}%
  \endgroup}%
  \paddedcell{\centering\SymbolFont\pairlines{c}{\ensuremath{\displaystyle #2}}{\ensuremath{\displaystyle #3}}}&%
  \paddedcell{\DefinitionFont\pairlines{l}%
    {\parbox[t]{\linewidth}{\raggedright #4}}%
    {\parbox[t]{\linewidth}{\raggedright #5}}}&%
  \paddedcell{\ExampleFont\renewcommand{\ExampleLineStretch}{1}%
    \exmath[#1]{%
    \pairstrut{\PairFirst}#6\\\pairstrut{\PairSecond}#7}}\\\hline
}
% Compact contexts use the existing palette and the table's example type.
\newcommand{\examplecontexts}[1]{%
  \begingroup
  \setlength{\fboxsep}{0.5mm}%
  \setlength{\fboxrule}{0.35pt}%
  \noindent\fcolorbox{coral}{coral!10}{%
    \begin{minipage}{\dimexpr\linewidth-2\fboxsep-2\fboxrule\relax}
      \ExampleFont\color{navy}%
      \textbf{Example contexts}\par
      #1\par
    \end{minipage}}\par
  \endgroup
  \vspace{0.7mm}%
}
% Optional contexts appear after the heading rule and before the table.
\newcommand{\sectionbox}[3][]{%
  {\HeadingFont\color{coral}#2\par}%
  \vspace{0.4mm}{\color{navy}\hrule height 0.5pt}\vspace{0.4mm}%
  \if\relax\detokenize{#1}\relax\else\examplecontexts{#1}\fi
  \begingroup
  \if\relax\detokenize{#1}\relax\else
    \renewcommand{\ExampleLineStretch}{1}%
  \fi
  \setlength{\SymbolWidth}{.24\linewidth}%
  \setlength{\MeaningWidth}{.32\linewidth}%
  \setlength{\tabcolsep}{0.7mm}%
  \renewcommand{\arraystretch}{1}%
  \arrayrulecolor{rulecolor}%
  \begin{tabularx}{\linewidth}{@{}>{\centering\arraybackslash\SymbolFont}p{\SymbolWidth}>{\raggedright\arraybackslash\DefinitionFont}p{\MeaningWidth}>{\raggedright\arraybackslash\ExampleFont}X@{}}
  \paddedcell{\centering\TableHeadFont Symbol}&\paddedcell{\TableHeadFont Meaning}&\paddedcell{\TableHeadFont Examples}\\\hline
  #3
  \end{tabularx}%
  \endgroup\par\vspace{0.8mm}%
}
% Dedicated two-column reference for letter names.
\newcommand{\letterrow}[2]{%
  \paddedcell{\centering\SymbolFont\ensuremath{#1}}&%
  \paddedcell{\DefinitionFont #2}\\\hline
}
\newcommand{\letterreference}[1]{%
  {\HeadingFont\color{coral}Letters and Greek symbols\par}%
  \vspace{0.4mm}{\color{navy}\hrule height 0.5pt}\vspace{0.4mm}%
  \begingroup
  \setlength{\tabcolsep}{0.7mm}%
  \renewcommand{\arraystretch}{1}%
  \arrayrulecolor{rulecolor}%
  \begin{tabularx}{\linewidth}{@{}>{\centering\arraybackslash\SymbolFont}p{.29\linewidth}>{\raggedright\arraybackslash\DefinitionFont}X@{}}
  \paddedcell{\centering\TableHeadFont Symbol}&\paddedcell{\TableHeadFont Name / description}\\\hline
  #1
  \end{tabularx}%
  \endgroup\par\vspace{0.8mm}%
}

\newcommand{\teachingheading}[1]{%
  {\HeadingFont\color{coral}#1\par}%
  \vspace{0.4mm}{\color{navy}\hrule height 0.5pt}\vspace{1.5mm}}
\newcommand{\teachinglabel}[1]{\vspace{1.8mm}{\DefinitionFont\bfseries\color{coral}#1\par}\vspace{0.7mm}}
\newcommand{\teachingmath}[1]{\par\vspace{1mm}{\ExampleFont\exmath{#1}\par}\vspace{1mm}}
\begin{document}
\begingroup

{\TitleFont\color{coral}MATHEMATICS\hfill\KickerFont POSTER 1\par}
\vspace{1.3mm}{\color{navy}\hrule height 0.8pt}\vspace{1.8mm}
{\KickerFont GENERAL NOTATION, SETS AND PROBABILITY\par}\vspace{3mm}

\begingroup
\setlength{\RowPadding}{1.8mm}
\sectionbox{General notation}{
\pairedrow{=}{\neq}{Equal}{Not equal}{2+3&=5;&4-4&=0}{2&\neq3;&-1&\neq1}
\symrow{\approx}{Approximately}{\exmath{\pi&\approx3.14\\\sqrt{2}&\approx1.41}}
\pairedrow{<}{>}{Less than}{Greater than}{2&<3;&-1&<0}{5&>4;&0&>-2}
\pairedrow{\leq}{\geq}{Less than or equal}{Greater than or equal}{2&\leq2;&-1&\leq3}{5&\geq4;&0&\geq0}
\symrow{\pm}{Plus or minus}{\exmath{3\pm1&=4\text{ or }2\\-2\pm3&=1\text{ or }-5}}
\symrow{|a|}{Absolute value}{\exmath[r@{}E@{\quad}r@{}E]{|-5|&=5;&|3|&=3}}
\symrow{\sqrt{a}}{Square root}{\exmath[r@{}E@{\quad}r@{}E]{\sqrt9&=3;&\sqrt{\frac14}&=\frac12}}
}

\endgroup
\vspace{3mm}
\noindent\begin{minipage}[t]{.49\linewidth}
\setlength{\RowPadding}{2.5mm}
\sectionbox{Sets and number systems}{
\symrow{\{a,b,c\}}{Set with three elements}{$\{2,4,6\};\quad\{-1,0,5\}$}
\symrow{\varnothing}{Empty set}{\exmath{\{1\}\cap\{2\}&=\varnothing\\\{x\in\mathbb R:x^2=-1\}&=\varnothing}}
\pairedrow{\in}{\notin}{Is element of}{Is not element of}{2&\in\{1,2\};&-1&\in\mathbb Z}{3&\notin\{1,2\};&\frac12&\notin\mathbb Z}
\symrow{\subseteq}{Is a subset of}{\exmath{\{1\}&\subseteq\{1,2\}\\\mathbb N&\subseteq\mathbb Z}}
\pairedrowwithcols{l}{\cup}{\cap}{Union}{Intersection}
  {\exmath{\{1\}\cup\{2\}&=\{1,2\}\\\varnothing\cup\{0\}&=\{0\}}}
  {\exmath{\{1,2\}\cap\{2\}&=\{2\}\\\{-1\}\cap\{4\}&=\varnothing}}
\pairedrow{\mathbb N}{\mathbb Z}{Natural numbers}{Integers}{3&\in\mathbb N;&8&\in\mathbb N}{-2&\in\mathbb Z;&0&\in\mathbb Z}
\pairedrow{\mathbb Q}{\mathbb R}{Rational numbers}{Real numbers}{\frac12&\in\mathbb Q;&-3&\in\mathbb Q}{\sqrt2&\in\mathbb R;&\frac14&\in\mathbb R}
\pairedrow{(a,b)}{[a,b]}{Open interval}{Closed interval}{1&\notin(1,3);&0&\in(-2,2)}{1&\in[1,3];&-2&\in[-2,0]}
}


\end{minipage}\hfill
\begin{minipage}[t]{.485\linewidth}
\setlength{\RowPadding}{2.5mm}
\sectionbox[{
\textbf{I}: fair die \qquad \textbf{II}: two independent fair coin tosses.\par
$\mathrm{H}$ = heads; $\mathrm{T}$ = tails.\par
\textbf{I}: $A=\{1,3,5\}$, $B=\{1,2\}$.\par
\textbf{II}: $A=\{\mathrm{HH}\}$, $B=\{\mathrm{HH},\mathrm{HT},\mathrm{TH}\}$.\par
}]{Probability and events}{


\symrow{\Omega}{Sample space: all possible outcomes}{\exmath[r@{\:}r@{}E]{\text{I:}&\Omega&=\{1,2,3,4,5,6\}\\\text{II:}&\Omega&=\{\mathrm{HH},\!\mathrm{HT},\!\mathrm{TH},\!\mathrm{TT}\}}}
\symrow{\omega}{An individual outcome}{\exmath[r@{\:}r@{}E@{}E]{\text{I:}&\omega&=3&{}\in\Omega\\\text{II:}&\omega&=\mathrm{HT}&{}\in\Omega}}
\symrow{A=\{1,3,5\}}{Event shown as a set of outcomes}{\exmath[r@{\:}r@{}E]{\text{I:}&A&=\{1,3,5\}\text{: odd}\\\text{II:}&A&=\{\mathrm{HH}\}\text{: two heads}}}
\symrow{A\subseteq\Omega}{An event is a subset of the sample space}{\exmath[r@{\:}r@{}E]{\text{I:}&\{1,3,5\}&\subseteq\Omega\\\text{II:}&A&\subseteq\Omega}}
\symrow{P(A)}{Probability of event $A$}{\exmath[r@{\:}r@{}E@{}E]{\text{I:}&P(A)&=\frac36&{}=\frac12\\\text{II:}&P(A)&=\frac14&}}
\symrow{A^c}{Complement of event $A$}{\exmath[r@{\:}r@{}E]{\text{I:}&A^c&=\{2,4,6\}\\\text{II:}&A^c&=\{\mathrm{HT},\mathrm{TH},\mathrm{TT}\}}}
\symrow{A\cup B}{Event $A$ or $B$ (or both)}{\exmath[r@{\:}r@{}E]{\text{I:}&A\cup B&=\{1,2,3,5\}\\\text{II:}&A\cup B&=\{\mathrm{HH},\!\mathrm{HT},\!\mathrm{TH}\}}}
\symrow{A\cap B}{Both events $A$ and $B$}{\exmath[r@{\:}r@{}E]{\text{I:}&A\cap B&=\{1\}\\\text{II:}&A\cap B&=\{\mathrm{HH}\}}}
}

\end{minipage}\par
\vfill
{\color{navy}\hrule height 0.5pt}\vspace{1.5mm}
{\FooterFont Symbols acquire meaning from their mathematical context.\par}
\endgroup
\newpage
\begingroup

{\TitleFont\color{coral}MATHEMATICS\hfill\KickerFont POSTER 2\par}
\vspace{1.3mm}{\color{navy}\hrule height 0.8pt}\vspace{1.8mm}
{\KickerFont FUNCTIONS, VARIABLES, SEQUENCES AND SUMS\par}\vspace{3mm}

\noindent\begin{minipage}[t]{.49\linewidth}
\setlength{\RowPadding}{1.85mm}
\sectionbox[{
\textbf{I}: $f:\mathbb R\to\mathbb R$, $f(x)=x^2$ \qquad 
\textbf{II}: $g:[0,\infty)\to\mathbb R$, $g(x)=\sqrt{x}$\par
}]{Function notation}{
\symrow{f}{Name of a function}{\exmath[r@{\:}r@{}E]{\text{I:}&f(x)&=x^2\\\text{II:}&g(x)&=\sqrt{x}}}
\symrow{f(x)}{Value of $f$ at input $x$}{\exmath[r@{\:}r@{}E@{}E]{\text{I:}&f(3)&=3^2&{}=9\\\text{II:}&g(4)&=\sqrt4&{}=2}}
\symrow{f:A\to B}{Function from set $A$ to set $B$}{\exmath[r@{\:}r@{}E@{}E]{\text{I:}&f&:\mathbb R&\to\mathbb R\\\text{II:}&g&:[0,\infty)&\to\mathbb R}}
\symrow{x\overset{f}{\mapsto} f(x)}{Input-to-output mapping}{\exmath[r@{\:}r@{}E@{\quad}r@{}E]{\text{I:}&x&\overset{f}{\mapsto} x^2;&3&\overset{f}{\mapsto}9\\\text{II:}&x&\overset{g}{\mapsto}\sqrt{x};&4&\overset{g}{\mapsto}2}}
\symrow{\operatorname{dom}(f)}{Domain of a function}{\exmath[r@{\:}r@{}E]{\text{I:}&\operatorname{dom}(f)&=\mathbb R\\\text{II:}&\operatorname{dom}(g)&=[0,\infty)}}
}
\vspace{2mm}\letterreference{
\letterrow{a,b,c,\ldots}{Common lowercase Latin letters}
\letterrow{A,B,C,\ldots}{Common uppercase Latin letters}
\letterrow{x,y,z}{Other common Latin letters}
\letterrow{\alpha,\ \beta,\ \gamma}{Alpha; beta; gamma}
\letterrow{\delta,\ \Delta}{Delta; uppercase delta}
\letterrow{\theta,\ \Theta}{Theta; uppercase theta}
\letterrow{\lambda,\ \Lambda}{Lambda; uppercase lambda}
\letterrow{\mu,\ \sigma}{Mu; sigma}
\letterrow{\omega,\ \Omega}{Omega; uppercase omega}
}
\vspace{2mm}{\HeadingFont\color{coral}Variable Examples\par}
\vspace{0.7mm}{\color{navy}\hrule height 0.5pt}\vspace{0.7mm}
{\DefinitionFont
A variable is a symbol that can represent different values.
Its allowed values and meaning depend on the context.\par
\vspace{1.3mm}
{\color{coral}\bfseries Possible values in separate cases}\par
{\ExampleFont Each choice below is a separate case.\par}
\vspace{0.7mm}
\textbf{Ordinary numbers:} $x$ can be $-2$, $0$, or $3.5$.\par
\vspace{0.7mm}
\textbf{Integers:} $n$ can be $1,2,3,\ldots$ when counting objects.\par
\vspace{0.7mm}
\textbf{Angles:} $\theta$ can be $30^\circ$, $60^\circ$, or $90^\circ$.\par
\vspace{0.7mm}
\textbf{Sets:} $A$ can be $\{1,2\}$ or $\{0,3,5\}$.\par
\vspace{0.7mm}
\textbf{Greek letters as numbers:} $\lambda$ can be $-1$, $0$, or $2$.\par
\vspace{1.3mm}
Letters do not have permanently assigned types. Uppercase
letters often represent sets, and Greek letters often represent
angles or numbers, but the context determines their meaning.\par
}


\end{minipage}\hfill
\begin{minipage}[t]{.485\linewidth}
\raggedright

\teachingheading{Lists, sequences and subscripts}
{\DefinitionFont
A \textbf{sequence} is an ordered list of terms. A \textbf{series}
is a sum of terms. Write sequences with parentheses and sets with braces.\par
\teachinglabel{Terms have positions}
The terms $a_1,a_2,a_3,\ldots,a_n$ occupy positions $1,2,3,\ldots,n$.
The notation $a_i$ means \textbf{the term at position $i$}. The subscript
$i$ is an \textbf{index}, not multiplication or an exponent.\par
\teachingmath{a_1&=\text{first term}\\a_2&=\text{second term}\\a_n&=\text{term at position }n}
\teachinglabel{One finite sequence, three forms}
An explicit list, its indexed terms and an indexed family can describe the same sequence:\par
\teachingmath{(2,4,6,8)&=(a_1,a_2,a_3,a_4)\\&=(a_i)_{i=1}^{4},\quad a_i=2i.}
In general, $(a_i)_{i=1}^{n}=(a_1,a_2,\ldots,a_n)$.\par
\teachinglabel{A rule generates the terms}
For $a_i=2i$, with $i=1,2,3,4$:\par
{\ExampleFont\exmath[r@{}E@{\qquad}r@{}E]{
a_1&=2(1)=2;&a_2&=2(2)=4;\\
a_3&=2(3)=6;&a_4&=2(4)=8.}\par}
The resulting sequence is $(2,4,6,8)$.\par
\teachinglabel{Order and repetitions matter}
$(2,4)\ne(4,2)$ and $(2,2,4)\ne(2,4)$.\par
For sets, order and repetition do not matter:\par
\teachingmath{\{2,4\}&=\{4,2\}=\{2,2,4\}.}
\teachinglabel{From a sequence to its sum}
The sum of the terms is $a_1+a_2+\cdots+a_n$.
Summation notation writes this more compactly.\par
}
\vspace{4mm}
\teachingheading{Summation notation}
{\DefinitionFont
\textbf{Read the parts}\par\vspace{1mm}
\begingroup
\setlength{\tabcolsep}{1mm}\arrayrulecolor{rulecolor}
\begin{tabularx}{\linewidth}{@{}>{\centering\arraybackslash}m{.25\linewidth}X@{}}
\centeredcell{\centering\SymbolFont $\displaystyle\sum_{i=1}^{n}a_i$}&
\centeredcell{\DefinitionFont $\sum$: summation symbol\\
$i$: index; $1$: lower bound\\
$n$: upper bound\\
$a_i$: expression being summed}\\\hline
\end{tabularx}\endgroup\par\vspace{1.5mm}
The index runs through consecutive integers from the lower bound to the upper bound, \textbf{inclusive}.\par
\teachingmath{\sum_{i=1}^{n}a_i&=a_1+a_2+\cdots+a_n.}
\teachinglabel{Sum the sequence above: $a_i=2i$}
\teachingmath{\sum_{i=1}^{4}a_i&=\sum_{i=1}^{4}2i\\[1mm]&=2+4+6+8=20.}
\teachinglabel{Other summands and bounds}
{\ExampleFont\exmath[r@{\quad}r@{}E]{
\text{Constant:}&\sum_{i=2}^{4}3&=3+3+3=9.\\[1.5mm]
\text{Index expression:}&\sum_{i=1}^{2}(i+3)&=4+5=9.\\[1.5mm]
\text{Squares:}&\sum_{k=1}^{3}k^2&=1+4+9=14.
}\par}
\teachinglabel{The index name can change}
\teachingmath{\sum_{i=1}^{4}i&=\sum_{k=1}^{4}k=1+2+3+4=10.}
Use the new index name consistently in the bounds and summand.\par
}

\end{minipage}\par
\vfill
{\color{navy}\hrule height 0.5pt}\vspace{1.5mm}
{\FooterFont Symbols acquire meaning from their mathematical context.\par}
\endgroup
\end{document}
